\chapter{自适应MCKF松耦合导航的工程实现与实验验证}
%纯包含关系的单个子标题删除，不是纯包含关系的单个子标题融合进标题
\section{松耦合自适应MCKF滤波的实现框架与协方差的设定}
本节在第二章松耦合误差模型与第三章 MCKF 理论基础上，给出完整的离散滤波实现方案\cite{izanloo2016kalman}。连续时间系统模型与观测模型分别写为
\begin{equation}
	\delta \dot{\bm{x}}_k = \bm{F}_{\text{SINS}}\,\delta\bm{x}_k + \bm{G}_{\text{SINS}}\,\bm{w}_k,
	\label{eq:ch4_continuous_process}
\end{equation}
\begin{equation}
	\delta\bm{z}^{\text{LC}}_k = \bm{H}^{\text{LC}}_k\,\delta\bm{x}_k + \bm{e}^{\text{LC}}_k.
	\label{eq:ch4_continuous_measure}
\end{equation}
其中，状态向量定义、系统矩阵与观测矩阵物理含义与第二章保持一致。

令相邻历元时间间隔为 $\Delta t_k=t_k-t_{k-1}$，则式(\ref{eq:ch4_continuous_process})可离散化为
\begin{equation}
	\delta\bm{x}_{k} = \bm{\Phi}_{k,k-1}\,\delta\bm{x}_{k-1} + \bm{\Gamma}_{k-1}\,\bm{w}_{k-1},
	\label{eq:ch4_discrete_process}
\end{equation}
其中
\begin{equation}
	\bm{\Phi}_{k,k-1}=\exp\!\left(\bm{F}_{\text{SINS},k-1}\Delta t_k\right),
\end{equation}
\begin{equation}
	\bm{\Gamma}_{k-1}=\int_0^{\Delta t_k}\exp\!\left(\bm{F}_{\text{SINS},k-1}\tau\right)\bm{G}_{\text{SINS},k-1}\,\mathrm{d}\tau.
	\label{eq:ch4_gamma_integral}
\end{equation}
在本文实现中，过程噪声由陀螺零偏与加计零偏驱动，且在单个积分步内近似时不变，因此可采用一阶积分近似
\begin{equation}
	\bm{\Gamma}_{k-1}\approx
	\begin{bmatrix}
		\bm{0}_{9\times6}\\
		\Delta t_k\bm{I}_{6\times6}
	\end{bmatrix},
	\label{eq:ch4_gamma_approx}
\end{equation}
该式对应工程实现中偏置子块噪声映射随步长线性缩放的形式。

基于离散模型，采用第三章给出的 MCKF 预测与迭代更新框架。预测阶段为
\begin{align}
	\hat{\delta\bm{x}}_{k|k-1} &= \bm{\Phi}_{k,k-1}\hat{\delta\bm{x}}_{k-1|k-1},\\
	\bm{P}_{k|k-1} &= \bm{\Phi}_{k,k-1}\bm{P}_{k-1|k-1}\bm{\Phi}_{k,k-1}^{\mathrm T}
	+ \bm{\Gamma}_{k-1}\bm{Q}_{k-1}\bm{\Gamma}_{k-1}^{\mathrm T}.
	\label{eq:ch4_predict_cov}
\end{align}
更新阶段采用固定点迭代形式
\begin{equation}
	\hat{\delta\bm{x}}_{k|k}^{(i+1)}=
	\hat{\delta\bm{x}}_{k|k-1}+\bm{K}_k^{(i)}
	\left(\delta\bm{z}_k^{\text{LC}}-\bm{H}^{\text{LC}}_k\hat{\delta\bm{x}}_{k|k}^{(i)}\right),
	\label{eq:ch4_mckf_iter}
\end{equation}
\begin{equation}
	\bm{K}_k^{(i)}=\tilde{\bm{P}}_{k|k-1}^{(i)}\left(\bm{H}^{\text{LC}}_k\right)^{\mathrm T}
	\left[\bm{H}^{\text{LC}}_k\tilde{\bm{P}}_{k|k-1}^{(i)}\left(\bm{H}^{\text{LC}}_k\right)^{\mathrm T}+\tilde{\bm{R}}_k^{(i)}\right]^{-1},
	\label{eq:ch4_mckf_gain}
\end{equation}
其中 $\tilde{\bm{P}}_{k|k-1}^{(i)}$ 与 $\tilde{\bm{R}}_k^{(i)}$ 由最大相关熵权重矩阵重标定得到，具体构造见第三章相关推导。

%\subsection{噪声协方差参数设定}
过程噪声模型沿用第二章状态方程设定，此处展开一阶高斯-马尔可夫过程，主要用于说明过程噪声参数的计算方式\cite{wang2018using}：
\begin{equation}
	\dot{\bm{b}}_g=-\frac{1}{\tau_g}\bm{b}_g+\bm{w}_g,
	\qquad
	\dot{\bm{b}}_a=-\frac{1}{\tau_a}\bm{b}_a+\bm{w}_a,
	\label{eq:ch4_markov_bias}
\end{equation}
其中 $\tau_g,\tau_a$ 为相关时间常数，可由 Allan 方差辨识结果给出；$\bm{w}_g,\bm{w}_a$ 为零均值白噪声。

设陀螺与加计稳态偏置标准差分别为 $\sigma_{bg},\sigma_{ba}$，则连续噪声谱密度可写为
\begin{equation}
	q_{bg}=\frac{2\sigma_{bg}^2}{\tau_g},
	\qquad
	q_{ba}=\frac{2\sigma_{ba}^2}{\tau_a}.
	\label{eq:ch4_psd_from_allan}
\end{equation}
进而过程噪声协方差采用如下块对角形式
\begin{equation}
	\bm{Q}_{k}=
	\operatorname{diag}
	\left(
	q_{bg}\bm{I}_{3\times3},
	q_{ba}\bm{I}_{3\times3}
	\right).
	\label{eq:ch4_Qk_def}
\end{equation}

松耦合模型的观测量为 GNSS 位置与速度解与 INS 解的差值。由于 INS 误差已在状态方程中显式建模，量测噪声项主要由 GNSS 传感器噪声决定，因此观测噪声协方差可设为
\begin{equation}
	\bm{R}^{\text{GNSS}}_k=
	\operatorname{diag}
	\left(
	\left(\sigma^{\text{GNSS}}_{r,E}\right)^2,\left(\sigma^{\text{GNSS}}_{r,N}\right)^2,\left(\sigma^{\text{GNSS}}_{r,U}\right)^2,
	\left(\sigma^{\text{GNSS}}_{v,E}\right)^2,\left(\sigma^{\text{GNSS}}_{v,N}\right)^2,\left(\sigma^{\text{GNSS}}_{v,U}\right)^2
	\right),
	\label{eq:ch4_Rk_def}
\end{equation}
其中 $\sigma^{\text{GNSS}}_{r,\cdot}$ 与 $\sigma^{\text{GNSS}}_{v,\cdot}$ 分别表示 GNSS 位置解与速度解在各方向上的标准差，可由接收机指标或静态数据统计获得。上述 $\bm{Q}_k$ 与 $\bm{R}^{\text{GNSS}}_k$ 的参数化形式可直接用于后续仿真实验章节的数据填充与对比分析。

%工程实现方法
%1 自适应带宽新息矩阵的非正定问题
% r_k_hat 采用了如下的方法进行正定保证 （以下方法表达为公式）
        % m = (m + m.T) / 2
        
        % # 2. 特征值修正
        % eigvals, eigvecs = np.linalg.eigh(m)
        % min_eig = np.min(eigvals)
        % if min_eig < eps:
        %     eigvals = np.maximum(eigvals, eps)
        %     m = eigvecs @ np.diag(eigvals) @ eigvecs.T
        
        % # 3. 可选：对角加载（双保险）
        % # R = R + eps * np.eye(R.shape[0])
    
        % return m

% 2对于不同更新频率gnss与ins的信息最大化利用问题
% 首先 一定应该保证时间戳的同步 在现实中 可以利用同步硬触发 ptp协议等多种方法保证 在此简略带过
% 常见的imu传感器的更新频率大概在100hz以上 而gnss单元普遍在1hz 为了充分利用imu数据 且保证算法的输出一直持续以满足某些实时场景的需求
% 工程设计上 当仅有imu数据进入时 执行imu数据的机械化编排以进行积分 然后输入mckf进行预测步骤 这里注意 仅当有imu数据进入时才能执行有效的预测 因为该预测模型的过程方程依赖imu的最新数据
% 与此同时 一直等待gnss数据 当有gnss数据的时间戳与imu数据的时间戳小于一定阈值时 认为他们是一组同步的数据 此时先按照imu单独处理步骤执行机械化编排 预测
%然后再进行后验更新 执行mckf的自适应带宽计算 和 标准更新步骤 
% 如此 最终输出的结果 经过模型中误差的补偿 在任意时间都是可用的


%3 注意到 上面lc模型的观测和过程其实都是线性的 因此 雅可比矩阵H和 F 其实就是观测矩阵和过程方程本身
%请你结合chap2中的插图 chap3的算法原理 上文提到的所有方法 画一个
%tikz
%的流程图 展示算法的全貌

\section{关键工程化处理方法}

本节给出算法在工程落地中的关键计算方法，重点说明在数值稳定性与异步传感器调度上的实现细节。以下方法均与前述松耦合线性误差模型和第三章 MCKF 迭代更新保持一致。

\subsection{自适应带宽新息协方差的正定化处理}

在自适应带宽更新中，常通过新息协方差匹配计算等效观测噪声矩阵 $\hat{\bm{R}}_k$。受有限样本与数值舍入影响，$\hat{\bm{R}}_k$ 可能出现非对称或非正定，进而导致 Cholesky 分解失败。为保证后续迭代稳定，本文采用“对称化 + 特征值修正 + 可选对角加载”的三步处理。

设待修正矩阵为 $\bm{M}$，先执行对称化：
\begin{equation}
	\bm{M}_s=\frac{1}{2}\left(\bm{M}+\bm{M}^{\mathrm T}\right).
	\label{eq:ch4_spd_sym}
\end{equation}
再对 $\bm{M}_s$ 做特征分解
\begin{equation}
	\bm{M}_s=\bm{U}\,\operatorname{diag}(\lambda_1,\ldots,\lambda_n)\,\bm{U}^{\mathrm T},
\end{equation}
并设定最小特征值阈值 $\varepsilon>0$，得到修正特征值
\begin{equation}
	\bar{\lambda}_i=\max(\lambda_i,\varepsilon),\quad i=1,\ldots,n.
\end{equation}
据此重构正定矩阵
\begin{equation}
	\bar{\bm{M}}=\bm{U}\,\operatorname{diag}(\bar{\lambda}_1,\ldots,\bar{\lambda}_n)\,\bm{U}^{\mathrm T}.
	\label{eq:ch4_spd_rebuild}
\end{equation}
在高噪声场景下可进一步使用对角加载
\begin{equation}
	\bar{\bm{M}}\leftarrow\bar{\bm{M}}+\varepsilon\bm{I}.
	\label{eq:ch4_spd_loading}
\end{equation}
工程实现中将 $\bar{\bm{M}}$ 作为 $\hat{\bm{R}}_k$ 的可分解近似，以保证 $\tilde{\bm{R}}_k$ 计算与固定点迭代可持续执行。

\subsection{异步采样条件下的双速率融合调度}

GNSS 与 IMU 典型更新频率存在显著差异（如 IMU 为 100 Hz 及以上，GNSS 为 1 Hz）。为最大化信息利用率并保证输出连续性，本文采用事件驱动的双速率调度策略。

设 IMU 时间戳为 $t_k^{\mathrm{imu}}$，GNSS 时间戳为 $t_j^{\mathrm{gnss}}$，同步判据为
\begin{equation}
	\left|t_k^{\mathrm{imu}}-t_j^{\mathrm{gnss}}\right|\leq \Delta t_{\mathrm{sync}},
	\label{eq:ch4_sync_rule}
\end{equation}
其中 $\Delta t_{\mathrm{sync}}$ 为工程容差阈值。系统按如下规则执行：
\begin{equation}
	\begin{aligned}
		\text{IMU-only event}:\quad
		\text{机械编排} \rightarrow \text{MCKF预测}
	\end{aligned}
\end{equation}
\begin{equation}
	\begin{aligned}
		\text{IMU-GNSS matched event}:\quad
		&\text{机械编排} \rightarrow \text{MCKF预测} \\
		&\rightarrow \text{自适应带宽更新} \rightarrow \text{MCKF校正}
	\end{aligned}
.
\end{equation}
该机制保证了两点：一是在无 GNSS 观测时仍可依赖 IMU 传播得到连续导航输出；二是在 GNSS 到达且时间对齐时及时完成后验修正，抑制纯惯导误差累积。

\subsection{线性误差模型下的简化计算}

由于本章松耦合误差模型在实现中为线性模型，过程雅可比与观测雅可比分别直接取系统矩阵与观测矩阵，即
\begin{equation}
	\bm{F}_k=\bm{F}_{\text{SINS},k},\qquad \bm{H}_k=\bm{H}^{\text{LC}}_k.
	\label{eq:ch4_linear_jacobian}
\end{equation}
如此一来，避免了公式雅可比矩阵求解带来的算力消耗或数值雅可比矩阵求解法带来的误差，简化了计算。

\subsection{误差状态的闭环不完全反馈策略}

在 GNSS/INS 松耦合误差状态滤波器中，估计得到的误差状态 $\hat{\delta\bm{x}}_k$ 需反馈回名义导航状态以实现误差补偿。为避免过度修正引发数值振荡与稳定性问题，本文采用闭环不完全反馈策略，即对不同误差分量设置独立反馈因子或直接采用单次反馈模式。

在具体实现中，每当完成一次 MCKF 更新，立即按如下形式进行误差反馈：
\begin{equation}
	\hat{\bm{x}}_{k|k} = \hat{\bm{x}}_{k|k-1} \oplus \hat{\delta\bm{x}}_{k|k},
\end{equation}
其中 $\oplus$ 表示适当的误差补偿算子，对位置、速度为矢量加法，对姿态为旋转更新。反馈后将误差状态重置为零，仅保留协方差矩阵 $\bm{P}_{k|k}$ 用于下一时刻预测，从而形成"估计-反馈-重置"的闭环不完全反馈机制。该策略既保证了误差得到及时抑制，又避免了误差状态与名义状态的重复计入，在工程实践中表现出良好的数值稳定性。

\subsection{算法整体流程图}

结合第二章模型结构与第三章 MCKF 原理，工程实现流程拆分为"主循环流程图"和"更新子模块流程图"，分别见图~\ref{fig:ch4_engineering_flow_main} 与图~\ref{fig:ch4_engineering_flow_correct}。


% 还可以继续凹
\begin{figure}[htbp]
	\centering
	\makebox[\textwidth][c]{\resizebox{0.94\textwidth}{!}{%
	\begin{tikzpicture}[
		node distance=6mm and 8mm,
		>=Latex,
		font=\dawu,
		data/.style={draw, ellipse, align=center, text width=34mm, minimum height=8mm, inner sep=1.5pt},
		sensor/.style={draw, rounded corners, align=center, text width=30mm, minimum height=8mm, inner sep=1.5pt},
		block/.style={draw, rounded corners, align=center, text width=38mm, minimum height=8mm, inner sep=2pt},
		decision/.style={draw, diamond, aspect=2.0, align=center, text width=26mm, inner sep=1.2pt, minimum width=22mm},
		flow/.style={-Latex, thick},
		loop/.style={-Latex, thick, dashed},
		group/.style={draw, rounded corners, dashed, inner sep=4pt}
	]
		\node[block] (init) {初始化\\$\hat{\delta\bm{x}}_0,\bm{P}_0,\bm{Q},\bm{R}^{\text{GNSS}}$};
		\node[decision, below=of init] (sync) {时间对齐?\\$|t_k-t_j|\leq\Delta t_{\text{sync}}$};
		\node[data, left=22mm and 8mm of sync] (imuData) {补偿惯性增量\\$\bm{u}^{c}_k=[\Delta\bm{\theta}^{c}_k,\Delta\bm{v}^{c}_k]^T$\\$\Delta\bm{\theta}^{c}_k=(\bm{\theta}^{m}_k-\hat{\bm{\varepsilon}}^b_k\Delta t_k)$\\$\phantom{\Delta\bm{\theta}^{c}_k=}-(\bm{\theta}^{m}_{k-1}-\hat{\bm{\varepsilon}}^b_{k-1}\Delta t_{k-1})$\\$\Delta\bm{v}^{c}_k=(\bm{v}^{m}_k-\hat{\nabla}^b_k\Delta t_k)$\\$\phantom{\Delta\bm{v}^{c}_k=}-(\bm{v}^{m}_{k-1}-\hat{\nabla}^b_{k-1}\Delta t_{k-1})$};
		\node[data, right=22mm and 8mm of sync] (gnssData) {GNSS观测向量\\$\bm{z}^{\mathrm{GNSS}}_k=[\bm{r}^{\mathrm{GNSS}}_k,\bm{v}^{\mathrm{GNSS}}_k]^T$};
		\node[sensor, above left=8mm and 6mm of imuData] (imuRaw) {IMU原始测量\\$\bm{y}^{\mathrm{imu},m}_k=[\bm{\theta}^{m}_k,\bm{v}^{m}_k]^T$};
		\node[sensor, above right=8mm and 6mm of gnssData] (gnssRaw) {GNSS原始测量\\$\bm{y}^{\mathrm{gnss},m}_k=[\tilde{\bm{r}}_k,\tilde{\bm{v}}_k]^T$};
		\node[block, below=of sync] (mech) {INS机械编排\\名义状态传播};
		\node[block, below=of mech] (pred) {Predict\\$\hat{\delta\bm{x}}^-_k=\bm{\Phi}_k\hat{\delta\bm{x}}^+_{k-1}$\\$\bm{P}^-_k=\bm{\Phi}_k\bm{P}^+_{k-1}\bm{\Phi}_k^T+\bm{\Gamma}_k\bm{Q}_k\bm{\Gamma}_k^T$};
		\node[decision, below=of pred] (hasGnss) {匹配GNSS观测?};
		\node[block, left=22mm of hasGnss] (predOut) {仅预测输出\\$\hat{\bm{x}}^-_k$};
		\node[block, right=22mm of hasGnss] (correct) {进入更新子模块\\见图~\ref{fig:ch4_engineering_flow_correct}};
		\node[block, below=11mm of hasGnss] (fb) {误差反馈\\$\hat{\bm{x}}_{k|k}=\hat{\bm{x}}_{k|k-1}\oplus\hat{\delta\bm{x}}_{k|k}$};
		\node[block, below=of fb] (out) {输出导航解};

		\draw[flow] (imuRaw.east) -- node[above, font=\scriptsize] {差分+偏置补偿} (imuData.west);
		\draw[flow] (gnssRaw.west) -- node[above, font=\scriptsize] {坐标/时间统一} (gnssData.east);
		\draw[flow] (imuData.east) -- node[above, font=\scriptsize] {$\bm{u}^{c}_k$} (sync.west);
		\draw[flow] (gnssData.west) -- node[above, font=\scriptsize] {$\bm{z}^{\mathrm{GNSS}}_k$} (sync.east);
		\draw[flow] (init) -- (sync);
		\draw[flow] (sync) -- (mech);
		\draw[flow] (mech) -- (pred);
		\draw[flow] (pred) -- (hasGnss);
		\draw[flow] (hasGnss) -- node[above, font=\scriptsize] {否} (predOut);
		\draw[flow] (hasGnss) -- node[above, font=\scriptsize] {是} (correct);
		\draw[flow] (predOut) |- (fb.west);
		\draw[flow] (correct) |- (fb.east);
		\draw[flow] (fb) -- (out);
		\draw[loop] (out.west) .. controls +(-22mm,-1mm) and +(-22mm,1mm) .. node[left, font=\scriptsize] {$k\leftarrow k+1$} (sync.west);

		\node[group, fit=(sync)(mech)(pred)(hasGnss)(correct)(fb), label={[font=\scriptsize]above:误差状态滤波器主模块}] {};
	\end{tikzpicture}%
	}}
	\caption{GNSS/INS松耦合自适应MCKF主循环流程（数据流与递归回路）}
	\label{fig:ch4_engineering_flow_main}
\end{figure}

\begin{figure}[htbp]
	\centering
	\makebox[\textwidth][c]{\resizebox{0.74\textwidth}{!}{%
	\begin{tikzpicture}[
		node distance=6mm and 8mm,
		>=Latex,
		font=\scriptsize,
		block/.style={draw, rounded corners, align=center, text width=46mm, minimum height=8mm, inner sep=2pt},
		decision/.style={draw, diamond, aspect=2.1, align=center, text width=26mm, inner sep=1.2pt, minimum width=22mm},
		flow/.style={-Latex, thick},
		loop/.style={-Latex, thick, dashed},
		group/.style={draw, rounded corners, dashed, inner sep=4pt}
	]
		\node[block] (innov) {更新-1: 自适应带宽\\$\hat{\bm{R}}_k\leftarrow\hat{\bm{P}}_{s,k}-\bm{H}_k\bm{P}^-_k\bm{H}_k^T$\\SPD修正: $\hat{\bm{R}}_k\succ0$};
		\node[block, below=of innov] (iterInit) {更新-2: 迭代初始化\\$\hat{\delta\bm{x}}_{k|k}^{(0)}\leftarrow\hat{\delta\bm{x}}_{k|k-1}$};
		\node[block, below=of iterInit] (iterStep) {MCKF固定点迭代\\$\bm{K}^{(i)}_k=\tilde{\bm{P}}^{(i)}\bm{H}_k^T(\bm{H}_k\tilde{\bm{P}}^{(i)}\bm{H}_k^T+\tilde{\bm{R}}_k^{(i)})^{-1}$\\$\hat{\delta\bm{x}}_{k|k}^{(i+1)}=\hat{\delta\bm{x}}^-_k+\bm{K}^{(i)}_k(\delta\bm{z}_k-\bm{H}_k\hat{\delta\bm{x}}_{k|k}^{(i)})$};
		\node[decision, below=of iterStep] (conv) {收敛判据\\$\|\Delta\hat{\delta\bm{x}}\|/\|\hat{\delta\bm{x}}\|<\varepsilon$?};
		\node[block, below=of conv] (post) {更新-3: 后验更新\\$\bm{P}^+_k=(\bm{I}-\bm{K}_k\bm{H}_k)\bm{P}^-_k$\\输出 $\hat{\delta\bm{x}}^+_k,\bm{P}^+_k$};

		\draw[flow] (innov) -- (iterInit);
		\draw[flow] (iterInit) -- (iterStep);
		\draw[flow] (iterStep) -- (conv);
		\draw[flow] (conv) -- node[right, font=\scriptsize] {是} (post);
		\draw[loop] (conv.east) .. controls +(12mm,0) and +(12mm,0) .. node[right, font=\scriptsize] {否, $i\leftarrow i+1$} (iterStep.east);

		\node[group, fit=(innov)(iterInit)(iterStep)(conv)(post), label={[font=\scriptsize]above:更新子模块}] {};
		\node[group, fit=(iterInit)(iterStep)(conv), label={[font=\scriptsize]above:MCKF 更新（固定点迭代）}] {};
	\end{tikzpicture}%
	}}
	\caption{更新 子模块内部流程（含自适应带宽与固定点迭代）}
	\label{fig:ch4_engineering_flow_correct}
\end{figure}


\section{仿真实验设计}

为验证前述全球导航卫星系统（GNSS）/惯性导航系统（INS）松耦合自适应最大相关熵卡尔曼滤波（MCKF）的有效性，本文在可控仿真环境下开展对比实验。实验设计遵循“统一轨迹、统一噪声模型、统一评价口径、仅替换滤波策略”的原则，保证结论可复现且具有可比性。

\subsection{仿真平台与实现边界}

本文采用 Aceinna 开源的 ``gnss-ins-sim'' 在 Python 环境生成带时间戳的 IMU 与 GNSS 观测数据，融合算法负责机械编排、状态预测与量测更新。数据读写、坐标变换及部分惯导基础工具参考 ``pyins'' 实现，松耦合自适应 MCKF 的误差状态建模、预测与迭代更新由本文独立完成\cite{gnssinssim,pyins}，并已开源发布\footnote{项目地址：\url{https://gitee.com/smallblackqaq/starlight}}。为保证结果可复现，本文对同一工况重复运行 5 次，并对各项指标取均值、标准差与显著性统计。

为避免平台差异对对比结果的影响，所有算法共用同一组仿真输入数据与初始条件，且统一采用第二章的 15 维松耦合误差状态定义与第三章的 MCKF 迭代更新框架。

\subsection{传感器误差模型与采样设置}

IMU 参数参考高精度惯导器件 Allan 方差指标进行设定，GNSS 观测噪声按仿真器字段 ``stdv'' 与 ``stdp'' 配置\cite{wang2018using,han2010quantization}。各项关键参数如表~\ref{tab:ch4_sim_sensor_params} 所示。

\begin{table}[htbp]
\centering
\caption{仿真实验传感器参数设置}
\label{tab:ch4_sim_sensor_params}
\begin{tabular}{ccc}
\toprule[1.5pt]
传感器 & 参数项 & 设定值 \\
\midrule[1pt]
IMU-陀螺 & 初始偏置 $\bm{b}_g$ & $[0,0,0]$ \\
IMU-陀螺 & 偏置漂移 $\bm{b}_{g,\text{drift}}$ & $[0.1,0.1,0.1]~(^{\circ}/\mathrm{h})$ \\
IMU-陀螺 & 相关时间常数 $\bm{b}_{g,\text{corr}}$ & $[100,100,100]~\mathrm{s}$ \\
IMU-陀螺 & 角随机游走 ARW & $[2.0\times10^{-3},2.0\times10^{-3},2.0\times10^{-3}]~(^{\circ}/\sqrt{\mathrm{h}})$ \\
\midrule[1pt]
IMU-加计 & 初始偏置 $\bm{b}_a$ & $[0,0,0]~\mathrm{m/s^2}$ \\
IMU-加计 & 偏置漂移 $\bm{b}_{a,\text{drift}}$ & $[3.6\times10^{-6},3.6\times10^{-6},3.6\times10^{-6}]~\mathrm{m/s^2}$ \\
IMU-加计 & 相关时间常数 $\bm{b}_{a,\text{corr}}$ & $[100,100,100]~\mathrm{s}$ \\
IMU-加计 & 速度随机游走 VRW & $[2.5\times10^{-5},2.5\times10^{-5},2.5\times10^{-5}]~\mathrm{m/s}/\sqrt{\mathrm{h}}$ \\
\midrule[1pt]
GNSS & 噪声参数 ``stdv'' & $[5.0,5.0,7.0]$ \\
GNSS & 噪声参数 ``stdp'' & $[0.05,0.05,0.05]$ \\
\bottomrule[1.5pt]
\end{tabular}
\end{table}

采样频率设置为 IMU 100 Hz、GNSS 1 Hz，对应前文提出的双速率融合调度。即在 IMU-only 事件执行机械编排与预测，在 GNSS 匹配事件执行“预测+更新”完整流程。

\subsection{轨迹与工况设计}

仿真轨迹采用城市道路风格长时程机动。初始位置设为纬度 $31.508183^{\circ}$、经度 $120.401989^{\circ}$、高度 $0~\mathrm{m}$；初始姿态为偏航 $183^{\circ}$、俯仰 $0^{\circ}$、横滚 $0^{\circ}$；初始速度为 $0~\mathrm{m/s}$。

按运动指令统计，整段轨迹时长为 $1410.3~\mathrm{s}$（约 $23.5~\mathrm{min}$），最大车体速度为 $20~\mathrm{m/s}$（约 $72~\mathrm{km/h}$）。轨迹由静止段、加减速段及多次转向机动交替组成，其中包含约 $200~\mathrm{s}$ 初始静止阶段，并在 $5/8/9.5/10/15/20~\mathrm{m/s}$ 多速度级之间切换；同时设置约 $\pm90^{\circ}$ 及更大幅度航向变化，以模拟城市道路转弯与姿态扰动。

该轨迹兼具“长时传播误差累积”与“机动段瞬态响应”两类特征，能够同时检验滤波器的稳态精度与动态鲁棒性。

\subsection{对比方案与评价指标}

为突出自适应 MCKF 在非理想噪声下的性能优势，本文采用以下两组算法进行同源数据对比\cite{izanloo2016kalman,xie2025adaptive}：
\begin{enumerate}
\item 标准扩展卡尔曼滤波（EKF）松耦合（固定 $\bm{Q},\bm{R}$）；
\item 自适应带宽最大相关熵卡尔曼滤波（MCKF）松耦合（本文方法）。
\end{enumerate}

%相关的统计数据 还做了CEP 同时每组结果计算PI

评价指标选取位置、速度与姿态三类误差的均方根值（RMSE）和峰值误差。以位置误差为例，其 RMSE 定义为
\begin{equation}
\mathrm{RMSE}_{r}=\sqrt{\frac{1}{N}\sum_{k=1}^{N}\|\hat{\bm{r}}_k-\bm{r}_k^{\text{ref}}\|_2^2},
\label{eq:ch4_rmse_pos}
\end{equation}
其中 $\hat{\bm{r}}_k$ 为估计位置，$\bm{r}_k^{\text{ref}}$ 为参考真值轨迹。速度与姿态指标采用同构定义。

对于分位统计指标，本文对位置、速度与姿态三类误差分别统计 50\% 与 95\% 分位数。位置通道先定义单次实验中第 $k$ 个时刻的水平误差为
\begin{equation}
e_{h,k}=\sqrt{\left(\hat r_{E,k}-r^{\text{ref}}_{E,k}\right)^2+\left(\hat r_{N,k}-r^{\text{ref}}_{N,k}\right)^2},
\label{eq:ch4_horizontal_error}
\end{equation}
其中 $E$ 和 $N$ 分别表示东向与北向分量。则单次实验的 CEP50 与 R95 可分别写为
\begin{equation}
\mathrm{CEP50}=Q_{0.50}\!\left(\{e_{h,k}\}_{k=1}^{N}\right),
\qquad
\mathrm{R95}=Q_{0.95}\!\left(\{e_{h,k}\}_{k=1}^{N}\right),
\label{eq:ch4_cep_r95}
\end{equation}
其中 $Q_{p}(\cdot)$ 表示 $p$ 分位数。速度与姿态通道采用同样的分位统计思想：分别对速度三维误差范数序列与姿态三维误差范数序列计算 50\% 与 95\% 分位数。若需要给出跨轮次统计值，则与 RMSE 一样，对各轮次对应分位指标再计算均值与标准差。

在单轮指标的基础上，本文对每组工况重复实验 5 次，因此表中的“$\pm$”表示不同实验轮次之间的均值与标准差，而不是单次轨迹内的离散波动。设第 $j$ 次实验得到的某一指标记为 $J_j$，重复次数为 $M$，则该指标的跨轮次均值与标准差分别定义为
\begin{equation}
\bar{J}=\frac{1}{M}\sum_{j=1}^{M}J_j,
\qquad
s_J=\sqrt{\frac{1}{M-1}\sum_{j=1}^{M}\left(J_j-\bar{J}\right)^2}.
\label{eq:ch4_mean_std}
\end{equation}
因此，若某一统计结果写为“$7.00\pm1.70$”，则表示 5 次重复实验得到的位置 RMSE 的均值为 $7.00~\mathrm{m}$、标准差为 $1.70~\mathrm{m}$。

% 此外，为考察抗异常能力，在 GNSS 观测中引入少量脉冲型粗差（幅值远高于标称噪声）并保持其发生时刻对各算法一致，比较不同方法在异常段的误差放大程度与收敛恢复时间。
%目前没搞这个哈 搞这个肯定要发散的

\subsection{实验流程与复现设置}

实验执行流程如下：
\begin{enumerate}
\item 依据表~\ref{tab:ch4_sim_sensor_params} 生成 IMU/GNSS 原始观测与真值轨迹；
\item 统一初值与协方差，分别运行 EKF 与自适应 MCKF，并对每组工况重复 5 次；
\item 依据式(\ref{eq:ch4_rmse_pos}) 及对应速度、姿态误差定义统计指标，同时记录位置/速度/姿态三通道的最大误差、50\% 分位数与 95\% 分位数；
\item 汇总时域误差曲线与统计表格，完成横向对比分析。
\end{enumerate}

% 为降低随机性影响，每组工况可在不同随机种子下重复多次并取均值统计。上述设计保证了后续实验结果能够直接对应“精度、稳定性、鲁棒性”三个维度，为第五章结果分析提供一致的数据基础。

%（为了降低随机性影响 充分发挥数据仿真平台的优势 实验采用不同的随机种子重复了20次）

\subsection{实验结果分析}

从典型单次运行曲线可以看出，两种方法在不同通道上的表现存在分化。图~\ref{fig:ch4_result_attitude} 与图~\ref{fig:ch4_result_heading} 显示，MCKF 在姿态与航向误差抑制方面整体更稳定，峰值尖峰与长尾漂移相对更弱。图~\ref{fig:ch4_result_velocity} 也表明 MCKF 在速度通道的波动幅度更小、异常时刻恢复更快。

但在位置通道，图~\ref{fig:ch4_result_position_3d} 与图~\ref{fig:ch4_result_position_h} 反映出 MCKF 并未稳定优于 EKF：在部分工况下其水平位置误差带更宽，说明当前参数设置下 MCKF 在 pose 指标上略逊于 EKF。综合图~\ref{fig:ch4_result_velocity_group} 的速度与轨迹结果，可以认为 MCKF 的主要收益集中在速度与姿态鲁棒性提升，而非位置均值误差的绝对最优。

\begin{figure}[htbp]
\centering
\begin{subfigure}[b]{0.49\textwidth}
\centering
\includegraphics[width=\textwidth]{experiment/typical_attitude_error_3d.png}
\caption{3D 姿态误差}
\label{fig:ch4_result_attitude}
\end{subfigure}
\begin{subfigure}[b]{0.49\textwidth}
\centering
\includegraphics[width=\textwidth]{experiment/typical_heading_error.png}
\caption{航向误差}
\label{fig:ch4_result_heading}
\end{subfigure}
\caption{典型单次运行的姿态误差结果}
\label{fig:ch4_result_attitude_group}
\end{figure}

\begin{figure}[htbp]
\centering
\begin{subfigure}[b]{0.49\textwidth}
\centering
\includegraphics[width=\textwidth]{experiment/typical_position_error_3d.png}
\caption{3D 位置误差}
\label{fig:ch4_result_position_3d}
\end{subfigure}
\begin{subfigure}[b]{0.49\textwidth}
\centering
\includegraphics[width=\textwidth]{experiment/typical_position_error_horizontal.png}
\caption{水平位置误差}
\label{fig:ch4_result_position_h}
\end{subfigure}
\caption{典型单次运行的位置误差结果}
\label{fig:ch4_result_position_group}
\end{figure}

\begin{figure}[htbp]
\centering
\begin{subfigure}[b]{0.49\textwidth}
\centering
\includegraphics[width=\textwidth]{experiment/typical_velocity_error_3d.png}
\caption{3D 速度误差}
\label{fig:ch4_result_velocity}
\end{subfigure}
\begin{subfigure}[b]{0.49\textwidth}
\centering
\includegraphics[width=\textwidth]{experiment/typical_trajectory_3d.png}
\caption{三维轨迹对比}
\label{fig:ch4_result_trajectory}
\end{subfigure}
\caption{典型单次运行的速度与轨迹结果}
\label{fig:ch4_result_velocity_group}
\end{figure}

五次重复实验的统计结果如表~\ref{tab:ch4_result_stats} 所示。位置通道上，MCKF 的位置 RMSE 为 $7.00\pm1.70~\mathrm{m}$，高于 EKF 的 $4.93\pm1.38~\mathrm{m}$；位置 CEP50 与位置 R95 分别为 $4.58\pm1.30~\mathrm{m}$、$13.12\pm3.43~\mathrm{m}$，也高于 EKF 的 $2.71\pm1.17~\mathrm{m}$、$6.63\pm2.90~\mathrm{m}$，说明 MCKF 在 pose 维度略逊。

相比之下，速度与姿态通道上 MCKF 优势显著：速度 RMSE 为 $0.247\pm0.058~\mathrm{m/s}$，显著低于 EKF 的 $0.429\pm0.475~\mathrm{m/s}$，且速度最大误差均值由 EKF 的 $4.09~\mathrm{m/s}$ 降至 $0.98~\mathrm{m/s}$。姿态 RMSE 方面，MCKF 为 $5.58\pm5.09^{\circ}$，明显低于 EKF 的 $22.95\pm18.02^{\circ}$；姿态 R95 也由 EKF 的 $62.94^{\circ}$ 降至 $15.16^{\circ}$。这表明自适应带宽机制在动态稳定性与抗异常扰动方面仍具有明显工程价值。

\begin{table}[htbp]
\centering
\caption{仿真实验结果统计（每组工况重复 5 次）}
\label{tab:ch4_result_stats}
\begin{tabular}{lcc}
\toprule[1.5pt]
指标 & MCKF & EKF \\
\midrule[1pt]
位置 RMSE / m & $7.00\pm1.70$ & $4.93\pm1.38$ \\
速度 RMSE / m/s & $0.247\pm0.058$ & $0.429\pm0.475$ \\
姿态 RMSE / $^{\circ}$ & $5.58\pm5.09$ & $22.95\pm18.02$ \\
最大位置误差 / m & $21.96\pm4.75$ & $17.08\pm9.08$ \\
最大速度误差 / m/s & $0.983\pm0.348$ & $4.093\pm4.318$ \\
最大姿态误差 / $^{\circ}$ & $21.18\pm16.73$ & $71.24\pm50.69$ \\
位置 CEP50 / m & $4.58\pm1.30$ & $2.71\pm1.17$ \\
位置 R95 / m & $13.12\pm3.43$ & $6.63\pm2.90$ \\
速度 CEP50 / m/s & $0.145\pm0.035$ & $0.086\pm0.048$ \\
速度 R95 / m/s & $0.489\pm0.106$ & $0.672\pm0.752$ \\
姿态 CEP50 / $^{\circ}$ & $0.461\pm0.162$ & $0.855\pm1.121$ \\
姿态 R95 / $^{\circ}$ & $15.16\pm14.42$ & $62.94\pm49.71$ \\
\bottomrule[1.5pt]
\end{tabular}
\end{table}

需要说明的是，位置 RMSE 的统计检验结果为 $p=0.0676$，在 $\alpha=0.05$ 水平下未达到显著性阈值，且均值层面 MCKF 劣于 EKF，这与“pose 略逊”的结论一致。另一方面，速度与姿态相关指标在均值与高分位统计上均显示出对 MCKF 的明显偏好，说明本章方法的主要收益体现在速度和姿态通道的鲁棒估计能力提升，而非位置误差最小化。

为便于观察跨轮次差异，表~\ref{tab:ch4_round_rmse} 列出了 5 轮实验中位置、速度与姿态 RMSE 的逐轮结果。可以看到，EKF 在第 3 轮出现了明显的速度与姿态失稳（速度 RMSE 为 $1.253~\mathrm{m/s}$、姿态 RMSE 为 $49.340^{\circ}$），显著抬升了其速度与姿态指标的标准差；而 MCKF 在第 4 轮姿态误差有阶段性上升（姿态 RMSE 为 $14.415^{\circ}$），对应图~\ref{fig:ch4_result_attitude_group} 中局部波动。该现象表明，两类方法均可能在特定工况下出现退化，但退化通道并不相同：EKF 主要表现为速度与姿态通道的突发失稳，MCKF 则更多体现为位置通道的持续偏差与个别姿态轮次波动。

此外，表~\ref{tab:ch4_pvalue_test} 给出了位置 RMSE 的统计检验结果。由于当前仅完成位置 RMSE 的检验，本章对速度与姿态“优势”的判断主要基于均值、标准差、分位数与典型时域曲线的一致性证据。后续可在配对实验设计下补充速度与姿态指标的同口径假设检验，以进一步增强结论的统计严谨性。

\begin{table}[htbp]
\centering
\caption{5轮实验逐轮 RMSE 对比}
\label{tab:ch4_round_rmse}
\begin{tabular}{ccccccc}
\toprule[1.5pt]
轮次 & \multicolumn{2}{c}{位置 RMSE / m} & \multicolumn{2}{c}{速度 RMSE / m/s} & \multicolumn{2}{c}{姿态 RMSE / $^{\circ}$} \\
\cmidrule(lr){2-3}\cmidrule(lr){4-5}\cmidrule(lr){6-7}
& MCKF & EKF & MCKF & EKF & MCKF & EKF \\
\midrule[1pt]
1 & 8.552 & 6.636 & 0.153 & 0.417 & 3.766 & 28.864 \\
2 & 6.397 & 4.406 & 0.296 & 0.191 & 1.691 & 23.127 \\
3 & 8.958 & 5.974 & 0.244 & 1.253 & 5.021 & 49.340 \\
4 & 6.171 & 4.437 & 0.245 & 0.106 & 14.415 & 2.234 \\
5 & 4.946 & 3.184 & 0.295 & 0.177 & 2.982 & 11.172 \\
\bottomrule[1.5pt]
\end{tabular}
\end{table}

\begin{table}[htbp]
\centering
\caption{位置 RMSE 显著性检验结果}
\label{tab:ch4_pvalue_test}
\begin{tabular}{lccc}
\toprule[1.5pt]
对比方法 & $p$ 值 & 判别阈值 $\alpha$ & 是否显著（$p<\alpha$） \\
\midrule[1pt]
MCKF vs EKF（位置 RMSE） & 0.0676 & 0.05 & 否 \\
\bottomrule[1.5pt]
\end{tabular}
\end{table}


%仿真结果分析
%xxxxxxxxxx

%此外 值得一提的是 作为